\chapter{\textbf{Elastoplasticity}}
\label{chap:plasticity}


\section*{Outline}
\addcontentsline{toc}{section}{\texorpdfstring{\hskip14mm}{} Outline}

This chapter treats the first genuinely nonlinear problem of the booklet.
It does so twice: first by programming an incremental algorithm by hand, so
that the structure of a nonlinear solver is explicit, and then by handing the
same problem to the standard \op{pasapas} procedure, which is what you will
actually use. Read both -- the second makes sense only once you have seen the
first.

The constitutive theory itself is only used here, not developed. For that,
the standard reference is Simo and Hughes~\cite{Simo:1998}: yield surfaces,
flow rules, hardening, the return-mapping algorithms that integrate them and
the consistent tangent operators that make Newton converge. Dhondt's
book~\cite{Dhondt:2004} covers the same material from the implementer's
side, in three dimensions and with thermomechanical coupling, and is the one
to read if you intend to write a constitutive routine of your own rather than
call one. Both are referred to again where the convergence output is
discussed in section~\ref{sec:plas:pasapas}.

\section{Plasticity -- programming of an incremental algorithm}
\label{sec:plas:incremental}

The example presented here is deliberately restricted to an elastoplastic
computation with a \emph{constant} stiffness matrix -- the initial elastic
stiffness is used at every iteration, so it is assembled and factorised once.
Convergence is therefore linear rather than quadratic, but each iteration is
cheap and the algorithm is robust. In this version the operator
\op{ecou[lement]} \dindex{ecou[lement]}{plastic flow, stress increment} does
the work of integrating the constitutive law, i.e. the radial return is not
programmed explicitly in \textsc{gibiane}.

\begin{castemcom}
We start from a \emph{square} domain \op{dom} -- note that this is not the
plate with a hole of the previous chapters, but a plain square built from
four straight edges -- on which the following loading is defined:

\begin{itemize}
\item imposed displacement
  \[ u_y |_{d34} = t\,\op{u1}, \qquad u_x |_{d23} = t\,\op{u2} \]
\item symmetry conditions on the two remaining edges, \op{d12} and
  \op{d41};
\item free surfaces elsewhere.
\end{itemize}

\noindent $t\in[0,1]$ is the fictitious time parameter and the loading is
applied in \op{ndt} steps of equal length. With \op{u1}${}=0$ the loading is
a simple tension in the $x$ direction; set it non-zero for a biaxial path.

The projection matrices associated with the blocked degrees of freedom are
\op{AA1}, \op{AA2}, \op{AA3} and their assembly \op{AA}, and the imposed
forces are stored for each step \op{n} in the force vector \op{FF . n} of
the table \op{FF}.
\end{castemcom}
\begin{castemlines}
\begin{verbatim}
* a plain square, independent of the
* plate of chapter 2.  It overwrites
* dom, d12 ... d41, so section 5.2
* rebuilds the plate before using it.
q1 = 0. 0.;  q2 = 1. 0.;
q3 = 1. 1.;  q4 = 0. 1.;
nc = 10;
d12 = droi nc q1 q2;
d23 = droi nc q2 q3;
d34 = droi nc q3 q4;
d41 = droi nc q4 q1;
dom = dall d12 d23 d34 d41;

u1 = 0.; u2 = 0.01; ndt = 8;

AA1 = (bloq uy d34 );  AA2 = (bloq ux d23);
AA3 = (bloq uy d12) et (bloq ux d41);
AA = AA1 et AA2 et AA3;

FF1 = depi AA1 u1; FF2 = depi AA2 u2;

FF = table;
  
repeter boucle (ndt + 1);
  n = &boucle - 1;  
  FF . n = ((flot n) / ndt ) * (FF1 et FF2);
fin boucle;

\end{verbatim}
\end{castemlines}
 
\begin{castemcom}
We first define the parameters of the constitutive behaviour:

\begin{itemize}
\item the Young modulus \op{EE} and the Poisson ratio \op{nu};
\item the yield stress \op{sigy} and the hardening modulus \op{H}.
\end{itemize}

\noindent An elastoplastic behaviour with \emph{kinematic} hardening is then
attached to the mesh \op{dom} with the \op{mode[le]} operator, and the
parameters of the model are attached to the model \op{mod} through
\op{mate[riaux]}.

The stiffness \op{KK} corresponding to the elastic part of the behaviour is
computed with \op{rigi[dite]} exactly as in the elastic case, so it depends
on \op{EE} and \op{nu} only. It is this matrix -- assembled once, before the
loop -- that makes the algorithm a constant-stiffness one.

\end{castemcom}
\begin{castemlines}
\begin{verbatim}

  EE = 2.e11;   nu = 0.3;
  sigy = 600.e6;
  H = 2000.e6;

  mod  = dom mode mecanique elastique
         isotrope plastique cinematique;
  mat = mate mod 'YOUN' EE 'NU' nu
              'H' H 'SIGY' sigy;

  KK = rigi mod mat;

\end{verbatim}
\end{castemlines}

\begin{castemcom}
\op{tol\_res} is the tolerance accepted on the equilibrium equation and
\op{it\_max} caps the number of equilibrium iterations per load step: these
two are the only numerical parameters of the computation.

The tables of plastic strains, internal variables, stresses and displacements
are initialised using \op{table} for the generalised list structure and
\op{zero} for the value at $t=0$. Note that \op{zero} needs the name of the
field type it is creating -- \texttt{'CONTRAIN'}, \texttt{'VARINTER'},
\texttt{'DEFINELA'} -- because the number and names of the internal variables
depend on the constitutive model declared above.
\end{castemcom}
\begin{castemlines}
\begin{verbatim}

  tol_res = 1.e-3;
  it_max  = 20;

  UU = table;    sig = table;
  epsp = table;  alpha = table;

  UU . 0 = manu 'CHPO' dom 2
                'UX' 0. 'UY' 0.;
  epsp . 0 = zero mod 'DEFINELA';
  alpha . 0 = zero mod 'VARINTER';
  sig . 0 = zero mod 'CONTRAIN';

\end{verbatim}
\end{castemlines}



\begin{castemcom}
The loop computing the load increments is named \op{it\_char}. It starts by
initialising the fields at step $n$ with the values converged at the previous
step $n-1$ -- the standard predictor of an incremental algorithm -- and then
computes a first elastic displacement increment \op{du} from the increment of
the applied load.
\end{castemcom}
\begin{castemlines}
\begin{verbatim}
repeter it_char ndt;
 n = &it_char; mess 'increment charge ' n;
    
 UU . n   = UU . (n - 1);
 epsp . n = epsp . (n - 1);
 sig . n   = sig . (n - 1);
 alpha . n = alpha . (n - 1);
 
 du = reso (KK et AA) (FF . n - FF . (n - 1 )); 
    
 UU . n = UU . n + du;
 deps = epsi mod du;

\end{verbatim}
\end{castemlines}

\pagebreak

\begin{castemcom}
Inside the load-increment loop \op{it\_char} a second loop \op{it\_equi} is
built. It contains the integration of the constitutive law -- the update of
the stresses and of the internal variables, performed by \op{ecou[lement]} --
followed by the verification of equilibrium.

\op{ecou} returns the new stress, the new internal variables and the
increment of plastic strain, given the state at the beginning of the step and
the total strain increment \op{deps}.
\end{castemcom}
\begin{castemlines}
\begin{verbatim}
 
repeter it_equi it_max;

nsig nalpha ndepsp = 
  ECOU mod (sig . n) (alpha . n) deps   mat;

sig . n   = nsig;
alpha . n = nalpha;
epsp . n  = (epsp . n) + ndepsp;
  
\end{verbatim}
\end{castemlines}

\begin{castemcom}
The verification of equilibrium starts with the computation of the residual
\op{RR}, defined as the difference between the imposed exterior forces
\op{FF . n} and the ``divergence'' of the stresses, i.e. the nodal forces
$[\mathbb{B}^T]\tens{\sigma}$ obtained with \op{bsigma}. The last term
removes the contribution of the imposed displacements from the nodal forces.

Between \op{si} and \op{finsi} we test whether the max norm of the residual
is smaller than the tolerance times the norm of the applied forces -- a
\emph{relative} criterion, which is what makes \op{tol\_res} independent of
the units of the problem -- and if so we \op{quit[ter]} the loop
\op{it\_equi}.

If the condition is not satisfied, a new displacement increment \op{du} is
computed to balance the residual and we return to the beginning of
\op{it\_equi} for another iteration: update of the internal variables and of
the stresses, and so on.
\end{castemcom}
\begin{castemlines}
\begin{verbatim}

RR = ((FF . n) - (bsigma mod sig . n)
              - (AA * (UU . n)));

no_res = n_force RR;
no_for = n_force ((FF . n) + (reac AA uu . n));

si ((maxi no_res) < (tol_res * (maxi no_for))); 
  quit it_equi;
finsi;

du = reso (KK et AA) RR;

UU . n = UU . n + du; deps = epsi mod du;

\end{verbatim}
\end{castemlines}


\begin{castemcom}
These two commands close the equilibrium loop \op{it\_equi} and the
load-increment loop \op{it\_char} respectively. Note that \op{fin} takes the
name of the loop it closes, which is what makes nested loops readable.
\end{castemcom}
\begin{castemlines}
\begin{verbatim}
fin it_equi;

fin it_char;
\end{verbatim}
\end{castemlines}

\begin{castemcom}
The commands \op{debp[rocedure]} and \op{finp[roc]} mark the beginning and
the end of the procedure \op{n\_force}. By extracting the components
\texttt{FX}, \texttt{FY} of the force field \op{ff} and combining them
algebraically, it computes the norm field
\[ \norm{\vect{f}} = \sqrt{f_x^2 + f_y^2}
   \quad \text{with} \quad
   \vect{f} = f_x \vect{e}_x + f_y \vect{e}_y \]

\medskip
\textbf{Place this procedure at the top of your file}, before the loops that
call it: \textsc{gibiane} is interpreted line by line, so a procedure must
have been read before it is invoked.
\end{castemcom}
\begin{castemlines}
\begin{verbatim}

debproc n_force ff*chpoint;
      
  lanorme =
   (((exco ff 'FX' 'SCAL')
     *(exco ff 'FX' 'SCAL')) +
    ((exco ff 'FY' 'SCAL')
     *(exco ff 'FY' 'SCAL'))) ** 0.5;

finproc lanorme; 

\end{verbatim}
\end{castemlines}

\section{Plasticity - computation using \texttt{pasapas} }
\label{sec:plas:pasapas}

\textbf{Note.} Section~5.1 overwrote \op{dom} and \op{d12}\ldots\op{d41}
with a plain square. Before running this section, rebuild the plate of
chapter~\ref{chap:meshing} -- or simply restore it with
\op{opti rest 'plate.dat'; rest;} if you saved it, as
chapter~\ref{chap:importexport} suggests.

\medskip
In this section we shall illustrate the usage of the standard \texttt{castem} command 
\texttt{pasapas} (fr. \emph{pas \`a pas}, step by step) command, which enables one to compute the solution of 
a nonlinear problem. The case discussed next will be that of an elastoplastic problem.
In the example we shall refer to the already defined geometry of the plate with a 
circular hole (see chapter~\ref{chap:meshing} and~\ref{chap:elasticity} ) and will directly 
start with the definition of the material, the boundary conditions. 

\subsubsection*{Model and loading history}

\begin{castemcom}
The \texttt{mode[l]} and the \texttt{mate[riaux]}  operators are used as
in the elastic case to define the material behaviour and its parameters.
In this case the model is a  elastoplastic model 
with  kinematic hardening. 
\end{castemcom}
\begin{castemlines}
\begin{verbatim}

   mod_p = dom mode mecanique
           elastique isotrope
                  plastique cinematique;
   mat_p   = mate mod_p 'YOUN' 2.e11  'NU' 0.3
                        'SIGY' 200.e6  'H' 2.e9;

\end{verbatim}
\end{castemlines}

As in the first version of the elastic computation in 
chapter~\ref{chap:elasticity}, we shall define an imposed 
displacement on the line $d_{34}$ of the plate. However 
if the value of the displacement was fixed in the elastic 
case, we shall define here a varying amplitude. 
The introduction of the imposed displacements will 
again be performed using a Lagrange multiplier technique. 


If the imposed displacement can be written as:
\begin{eqnarray}
  \vect{u}^D(x,t) &=& a(t) \, \vect{u}_0(x)
                      \qquad x \in \partial \Omega^d
\nonumber   
\end{eqnarray}
Then the projection matrix $[\mathbb{A}]$ will be the same with 
the one defined in the elastic case and will 
not vary in time. We shall just have to introduce 
the displacement amplitude in the program $a(t)$. 

\begin{castemcom}
We start again by defining the projection matrices and 
assembles them into only one element: $[\mathbb{A}]$. 

Next we define the associated nodal force vector 
$[\mathbb{U}_0^D]$ of the imposed displacement distribution 
$\vect{u}_0$, denoted as \texttt{ud34} in the program.  
 
The amplitude $a(t)$ is defined by its graph. In \texttt{castem}
this corresponds to a \texttt{evolution} object assembled from 
the list of time instants \texttt{tt} and the corresponding 
amplitudes \texttt{a\_ud34}.     

The association of the amplitude \texttt{ev\_ud34} with the 
corresponding nodal force vector \texttt{ud34} is done using 
the \texttt{char[gement]} (fr. \emph{chargement}) operator. The string 
\texttt{'DIMP'} announces that the loading is an 
imposed displacement. The keywords for different loading types
is specified in the help of the  \texttt{char[gement]}  
and the  \texttt{pasapas} operators.  

\end{castemcom}
\begin{castemlines}
\begin{verbatim}

  Ad12 = bloq uy d12;
  Ad41 = bloq ux d41;
  
  Ad34 = bloq uy d34;
  ud34  = depi Ad34 0.001;
  
  AA = Ad12 et Ad41 et Ad34;

   tt     = prog 0. 5. 10. 15. 20.;  
   a_ud34 = prog 0. 1. 0. -1. 0.;
   
   ev_ud34 = evol manu 'TEMPS' tt 'DIMP' a_ud34;
   dess ev_ud34;
   
   ch_ud34 = char 'DIMP' ud34 ev_ud34; 
 
 
\end{verbatim}
\end{castemlines}

\subsubsection*{The computation using \texttt{pasapas}}

\begin{castemcom}
The nonlinear computation is performed using the \texttt{pasapas} 
operator.
Its input variables are organized in a large \texttt{table}, 
that means a generalized list, where the different indexes 
specify the objects:
\begin{itemize}
  \item \texttt{'MODELE'} - the model
  \item \texttt{'CARACTERISTIQUES'} - the material parameters
  \item \texttt{'BLOCAGES\_MECANIQUES'} - the projection matrices for imposed displacement or temperature  
  \item \texttt{'CHARGEMENT'} - the loads, expressed as histories of nodal forces, constructed with 
      the \texttt{char[gement]} operator 
  \item \texttt{'TEMPS\_CALCULES'} - the \texttt{list} of time instant where equilibrium has to be checked
\end{itemize} 

\end{castemcom}
\begin{castemlines}
\begin{verbatim}

   tabexp = table;
      
   tabexp . 'MODELE' = mod_p;
   tabexp . 'CARACTERISTIQUES'  =   mat_p;

   tabexp . 'BLOCAGES_MECANIQUES' = AA;
   tabexp . 'CHARGEMENT'          = ch_ud34; 
 
   tabexp . 'TEMPS_CALCULES' =
                 prog 0. pas 1. 20.;

   pasapas tabexp; 


\end{verbatim}
\end{castemlines}

If the information provided for \texttt{pasapas} is consistent 
then the computation starts and during the performance a series 
of parameters are displayed for its monitoring. The output of the 
example will be discussed next. 

For each step number (\texttt{Numero du pas}), or the corresponding time
instant (\texttt{Indice d evolution}), one line is printed per iteration of
the inner loop that seeks plastic admissibility and global equilibrium.
Reading those six columns is the difference between watching a computation
and understanding it.

\cindex{convergence monitoring}\sindex{Nplas}\sindex{Critere}%
\sindex{Deps.max}\sindex{Eps.max}\sindex{Crit.flex}%
\begin{center}
\fcolorbox{black}{Gray}{%
\begin{tabular}{@{\hspace{2mm}}>{\ttfamily}l@{\hspace{5mm}}p{86mm}@{\hspace{2mm}}}
\multicolumn{2}{c}{\textbf{The \texttt{pasapas} convergence columns}} \\[2.mm]
Iter      & the iteration number within the current step \\[1.mm]
Nplas     & the \emph{number of integration points} whose internal variable
            changed by more than $10^{-10}$ during this iteration, i.e. how
            many points are currently flowing plastically \\[1.mm]
Critere   & the convergence criterion on forces -- the quantity that must
            fall below \texttt{'PRECISION'} \\[1.mm]
Deps.max  & the largest \emph{increment} of cumulated plastic strain produced
            by this iteration, $\max |\Delta p|$ \\[1.mm]
Eps.max   & the largest \emph{cumulated} plastic strain reached anywhere in
            the structure so far, $\max p$ \\[1.mm]
Crit.flex & the same criterion as \texttt{Critere} but formed on the
            \emph{moment} components -- \emph{flexion}, bending. Non-zero
            only for elements with rotational degrees of freedom: beams,
            plates, shells \\
\end{tabular}}
\end{center}

\noindent The two criteria are relative, which is what makes
\texttt{'PRECISION'} independent of your unit system. Writing $R$ for the
residual nodal force field, $U^{D}$ for the imposed displacements and $A$ for
the projection matrix of the blocked degrees of freedom,
\begin{equation}
  \texttt{Critere} \;=\;
  \frac{\max\Bigl(\;\max_{\text{force dof}} |R| \;,\;
                   \max |(U-U^{D})A| \;\Bigr)}{F_{ref}}
\label{eq:plas:critere}
\end{equation}
\begin{equation}
  \texttt{Crit.flex} \;=\;
  \frac{\max_{\text{moment dof}} |R|}{M_{ref}}
\label{eq:plas:critflex}
\end{equation}

\noindent Three things in~\eqref{eq:plas:critere} are worth noticing. The
numerator is a \emph{maximum}, not a Euclidean norm, so one badly converged
node is enough to hold up the whole step. It takes the larger of the force
residual and the violation of the imposed displacements, so a boundary
condition that is not being satisfied shows up in exactly the same number.
And the denominators $F_{ref}$, $M_{ref}$ are reference magnitudes built
internally from the applied loads and the reactions; when you set
\texttt{'FTOL'} or \texttt{'MTOL'} they are replaced by 1, the criterion
becomes an \emph{absolute} residual, and the column headings change to
\texttt{Fresidu} and \texttt{Mresidu} accordingly -- a useful signal that
you are no longer reading a relative number.

Convergence of the step requires \emph{both} criteria to fall below their
tolerances: \texttt{'PRECISION'}, default $10^{-4}$, and its moment
counterpart. The strain columns carry no tolerance; they are there to be
read by you.

\subsection*{What the columns tell you in practice}

\begin{itemize}
\item \texttt{Critere} should fall by roughly an order of magnitude per
  iteration. Falling and then stalling at some small value usually means the
  step is too large for the constitutive integration;
\item \texttt{Nplas} tells you how much of the structure is yielding. It
  rising steadily from step to step is the plastic zone spreading, which is
  physics. It oscillating wildly \emph{within} one step -- points yielding
  and unloading from one iteration to the next -- is the classic sign of a
  load increment too large to resolve the elastic-plastic boundary;
\item \texttt{Deps.max} is the quantity to watch against the step size. A
  plastic strain increment of a few percent in one step is too much for
  most constitutive laws; reduce \texttt{'TEMPS\_CALCULES'} until it is
  small, and note that \op{pasapas} will subdivide the step itself if it has
  to;
\item \texttt{Eps.max} is the accumulated plastic strain. Compare it with the
  ductility of the material: if it reaches tens of percent, the small-strain
  assumption of this chapter has stopped being valid and you belong in
  chapter~\ref{chap:largestrain};
\item \texttt{Crit.flex} identically zero simply means your model has no
  rotational degrees of freedom. On a shell model it not converging while
  \texttt{Critere} does is the signature of a bending problem -- usually too
  coarse a mesh through the thickness, or a membrane locking issue.
\end{itemize}

\noindent The algorithms behind all of this -- return mapping, the consistent
tangent operator, and why the choice of integration scheme governs the
convergence rate you see in the \texttt{Critere} column -- are the subject of
Simo and Hughes~\cite{Simo:1998}, which is the standard reference on
computational plasticity. Dhondt~\cite{Dhondt:2004} covers the same ground
from the point of view of someone implementing it in a three-dimensional
code, and is the more practical of the two if you intend to write your own
constitutive routine.

\begin{codebox}\footnotesize
\begin{verbatim}

 
 ------------------ DEBUT DE LA PROCEDURE PASAPAS ------------------
 
 
 Calcul MECANIQUE
  *** PLASTICITE ***
  *** CONVERGENCE FORCEE   ***
 
 
  Numero du pas :     1     Indice d evolution :   3.0000
  Iter   Nplas     Critere        Deps.max       Eps.max      Crit.flex
 Pas d increment de charge, initialisation calculee avec le temps
 Initialisation a partir de la solution precedente Coeff  0.0000
     1      35    0.50000         8.40624E-06     2.14937E-03     0.50000
     2      34    1.17475E-04     8.79403E-06     2.15165E-03     1.17475E-04
     3      34    8.47313E-05     1.14535E-05     2.15337E-03     8.47313E-05
  ****** CONVERGENCE A L ITERATION     3
   
  Numero du pas :     2     Indice d evolution :   4.0000
 
  Iter   Nplas     Critere        Deps.max       Eps.max      Crit.flex
     1     248    0.11013         5.67786E-04     2.72116E-03     0.11013
     2     255    1.42363E-03     7.69312E-04     2.92268E-03     1.42363E-03
     3     253    9.79151E-04     9.04782E-04     3.05815E-03     9.79151E-04
 act3 : reduction at 3 dimensions
     4     250    1.22652E-03     1.18403E-03     3.33740E-03     1.22652E-03
     5     251    2.00070E-04     1.21039E-03     3.36376E-03     2.00070E-04
     6     243    2.03488E-04     1.28299E-03     3.43636E-03     2.03488E-04
     7     240    9.73268E-05     1.29365E-03     3.44703E-03     9.73268E-05
  ****** CONVERGENCE A L ITERATION     7
   
   ....
   
   ------------------- FIN DE LA PROCEDURE PASAPAS -------------------
\end{verbatim}
\end{codebox}

\noindent The transcript above comes from a different run than the one set up
here, so do not expect your own step numbers and residuals to match it line
by line. What matters is the shape: a decreasing \texttt{Critere} and a
\texttt{CONVERGENCE A L ITERATION} line closing each step. A step that
reaches the iteration limit without that line, or a \texttt{Critere} that
stops decreasing, means the step did not converge -- reduce the step size in
\texttt{'TEMPS\_CALCULES'} before believing the results.


\subsubsection*{Analysis of the results}

After the completion of the \texttt{pasapas} computation 
a series of results can be analysed. 

\begin{castemcom}
\texttt{list}-ing the contents of \texttt{tabexp} 
presents the complete list of the entries of the table,
ranging from computational options and settings to the
fields of the solution. Some of the important entries 
are:
\begin{itemize}

 \item \texttt{deplacements} - the displacement fields 
 \item \texttt{contraintes} - the stress fields 
 \item \texttt{variables\_internes} - the  fields of the internal variables of the 
   constitutive model, its component fields are described in the manual pages 
   of \texttt{mode[le]} and \texttt{mate[riaux]}, as well as in documents 
   like~\cite{Pasquet:1998}  
 \item \texttt{deformations\_inelastiques} - the fields of the inelastic strains 

\end{itemize}
All these entries are themselves tables indexed by the number of the 
computed step. The correspondence between steps and time instants 
can easily be recovered by listing the table of time \texttt{tabexp . temps}.  

The global strain field is not saved in the procedure, but can 
easily be computed with \texttt{epsi[lon]}.

The computed fields can be displayed  
using \texttt{trac[er]}, in the usual way, that means 
appending the mesh for the nodal fields, \texttt{chpoint},  and the 
model for the fields on elements (\texttt{chamelem}, defined for example at the Gauss points).   

\index{extr[aire]@ \texttt{extr[aire]} extract, ! \texttt{'COMP'}, components    } 

\end{castemcom}
\begin{castemlines}
\begin{verbatim}

   list tabexp;

   list tabexp . deplacements ;
   list tabexp . contraintes ;
   list tabexp . variables_internes ;

   list tabexp . temps ;

   trace mod_p (tabexp . contraintes . 15);
   trace dom tabexp . deplacements . 4 ;

   list (extr
     (tabexp . variables_internes . 5) 'COMP');


\end{verbatim}
\end{castemlines}
     
\begin{castemcom}
If one wants to track the history  of a field at a certain 
point, one can easily construct the corresponding 
\texttt{evolution} object.

The \texttt{prog[ramer]} command initiates and then
constructs by appending the list of values for the time 
and field evolutions.   

The \texttt{repe[ter], fin} command permits to construct 
the loop over all time steps. The user defined  string  
\texttt{loop} is only a flag indicating the start and the 
end of the loop. This functionality is of importance for 
intricate loops. \texttt{\&loop} is the iterator of the 
loop \texttt{loop} and varies from 1 to 
\texttt{(dime tabexp . deplacements)}. However as 
the lists in \texttt{tabexp} start with the initial 
step set at \texttt{0} we insert the new iterator 
\texttt{ii}.

The use of \texttt{dime[nsion]} permits to extract  automatically  
the number of time instants.  

Finally let us note the usage of \texttt{extr[aire]}:
\begin{itemize}

   \item for nodal fields, \texttt{chpoint}'s , like the 
   displacement or the temperature field, one needs to indicate 
   the component and the nodal point of extraction, here 
   \texttt{'UX'} and \texttt{a3} respectively.    
   
   \item  for fields defined on elements, \texttt{chamelem}'s , 
   like the stress or the strain field, one needs to indicate 
   the component and the numbers of the zone, the element and of the Gauss point of 
   extraction, here \texttt{'SMXY'} and \texttt{1 30 2} respectively.    
   
   In different cases the models are automatically partitioned in zones
   that indicate for example different types of elements or 
   material behaviour.  

\end{itemize} 
 
\index{evol[ution]@ \texttt{evol[ution]} ! \texttt{'MANU'}, from lists } 
\index{extr[aire]@ \texttt{extr[aire]} extract, ! \texttt{chpoint}, nodal field    } 
\index{extr[aire]@ \texttt{extr[aire]} extract, ! \texttt{chamelem}, field on elements    } 
\dindex{repe[ter]}{do, start}
 
\end{castemcom}
\begin{castemlines}
\begin{verbatim}

   thetime = prog;
   uuy = prog;
   sigyy = prog;
   sigxy = prog;
   
   repeter loop (dime tabexp . deplacements);
   ii = &loop - 1;
   
   thetime = thetime et
             (prog (tabexp . temps . ii));
   
   uuy = uuy et (prog (extr
     (tabexp . deplacements . ii) 'UY' a3));
   
   sigyy = sigyy et 
      (prog (extr
        (tabexp . contraintes . ii)
        'SMYY' 1 30 2));
   sigxy = sigxy et 
      (prog (extr
        (tabexp . contraintes . ii)
        'SMXY' 1 30 2));
   
   fin loop;
   
\end{verbatim}
\end{castemlines}


\begin{castemcom}

After ending the loop we have retrieved the lists of time and 
of $u_y$, $\sigma_{xy}$, $\sigma_{yy}$ in the corresponding 
points. The lists can now be combined into \texttt{evolution}'s
which can be plotted using the \texttt{dess[iner]} command. 
 
\end{castemcom}
\begin{castemlines}
\begin{verbatim}

   ev_uy = evol rouge manu
       'time' thetime 'uy' uuy;
   ev_sigyy = evol jaune manu
       'time' thetime 'stress_yy' sigyy;
   ev_sigxy = evol vert manu
       'time' thetime 'stress_xy' sigxy;
   
   dess ev_uy;
   dess (ev_sigyy et ev_sigxy);

\end{verbatim}
\end{castemlines}

\section*{Summary}
\addcontentsline{toc}{section}{\texorpdfstring{\hskip14mm}{} Summary}

A nonlinear computation is two nested loops: an outer loop over load
increments, and an inner loop that iterates until the residual -- applied
forces minus $[B^{T}]\sigma$ -- is small compared with the applied forces.
Writing it out in \textsc{gibiane}, as in section~\ref{sec:plas:incremental},
takes about forty lines and is worth doing once.

After that, use \op{pasapas}. It performs the same two loops, with a far
better convergence strategy and automatic step subdivision, and it is driven
entirely by one \op{table}: you fill \texttt{'MODELE'},
\texttt{'CARACTERISTIQUES'}, \texttt{'BLOCAGES\_MECANIQUES'},
\texttt{'CHARGEMENT'} and \texttt{'TEMPS\_CALCULES'}, and it fills
\texttt{'DEPLACEMENTS'}, \texttt{'CONTRAINTES'} and
\texttt{'VARIABLES\_INTERNES'} back into the same table, each a sub-table
indexed by step. That input/output table is the single idea to take away: it
reappears unchanged for large strain in chapter~\ref{chap:largestrain}.

\section*{Exercises}
\addcontentsline{toc}{section}{\texorpdfstring{\hskip14mm}{} Exercises}

\begin{enumerate}
\item \textbf{Hand-written against \op{pasapas}.} Run the incremental
  algorithm of section~\ref{sec:plas:incremental} and the \op{pasapas} computation of
  section~\ref{sec:plas:pasapas} on the same square domain, with the same
  material and the same loading, and compare the stress--strain curves at a
  Gauss point. They should coincide to within the equilibrium tolerance.

\item \textbf{Cost of the tolerance.} Rerun the hand-written algorithm with
  \op{tol\_res} equal to $10^{-2}$, $10^{-3}$ and $10^{-5}$, counting the
  total number of equilibrium iterations. Plot the error on the final stress
  against the number of iterations.

\item \textbf{Kinematic versus isotropic hardening.} The example uses
  kinematic hardening. Change the model to isotropic hardening and apply a
  full load--unload--reload cycle. Which of the two reproduces the
  Bauschinger effect, and how do you see it on the stress--strain curve?

\item \textbf{Constant versus tangent stiffness.} The algorithm here
  re-uses the elastic stiffness at every iteration. Estimate how many
  iterations per step it needs, and discuss what a tangent-stiffness update
  would change in cost per iteration and in number of iterations.
\end{enumerate}
